\documentclass[UTF8,a4paper,11pt]{ctexart}
\usepackage{amsmath,amssymb}
\usepackage{geometry}
\usepackage{booktabs}
\usepackage{longtable}
\usepackage{hyperref}
\geometry{margin=2.2cm}
\title{子报告 1：UEP/ICCP 静态或慢变电场计算说明}
\author{}
\date{}
\begin{document}
\maketitle

\section{本源项计算目标}
本子报告解释主报告表 5 中“UEP/ICCP 静态或慢变电场”这一行如何计算。这里的 UEP/ICCP 不是把 UEP 和 ICCP 当成两个独立源相加，而是把腐蚀、防腐和外加电流共同形成的等效电流偶极矩记为 $P_0$，再计算它在海水中的静态或慢变电场。

\section{物理图像}
可以把水下目标看成一个“有电流从某些位置流出、又从另一些位置流回”的整体。远处探测器看不到所有局部细节，只能看到一个等效的电流偶极。电流偶极越大，远处电场越强；海水电导率越高，电场越容易被导走，测到的电场越弱；距离越远，电场按 $R^{-3}$ 快速下降。

\section{使用公式}
\[
P_0=P_{\mathrm{corr}}+P_{\mathrm{SACP}}+P_{\mathrm{ICCP}}
\]
\[
E_{\mathrm{static}}(R)=\frac{C_\theta P_0}{4\pi\sigma R^3}
\]

\begin{longtable}{p{2.6cm}p{3.2cm}p{2.2cm}p{5.5cm}}
\caption{公式变量说明}\\
\toprule
符号 & 名称 & 单位 & 来源和含义\\
\midrule
$P_0$ & 总等效电流偶极矩 & A$\cdot$m & 参数包络输入；代表腐蚀、防腐、ICCP 机制合成后的远场等效强度。\\
$P_{\mathrm{corr}}$ & 腐蚀相关偶极矩 & A$\cdot$m & 源项分解概念；主表不单独扫描。\\
$P_{\mathrm{SACP}}$ & 牺牲阳极防腐偶极矩 & A$\cdot$m & 源项分解概念；主表不单独扫描。\\
$P_{\mathrm{ICCP}}$ & 外加电流防腐偶极矩 & A$\cdot$m & 源项分解概念；主表不单独扫描。\\
$E_{\mathrm{static}}$ & 静态或慢变电场强度 & V/m & 需要输出的结果。\\
$C_\theta$ & 方向因子 & 无量纲 & 参数包络输入；横向取 1，轴向取 2，用来覆盖观测方向差异。\\
$\sigma$ & 水体电导率 & S/m & 参数包络输入；海水越导电，电场越小。\\
$R$ & 探测器到等效源中心距离 & m & 主表距离列；取 100、300、500、1000 m。\\
$4\pi$ & 球面几何常数 & 无量纲 & 来自三维空间中点源/偶极远场解析形式。\\
\bottomrule
\end{longtable}

\section{参数扫描}
\[
P_0=\{30,100,300,1000,3000\}\ \mathrm{A\cdot m}
\]
\[
\sigma=\{3,4,6\}\ \mathrm{S/m},\qquad C_\theta=\{1,2\}
\]
\[
R=\{100,300,500,1000\}\ \mathrm{m}
\]
参考场景取 $P_0=300\ \mathrm{A\cdot m}$、$\sigma=4\ \mathrm{S/m}$、$C_\theta=1$。

\section{Python 计算对应关系}
核心函数位于：
\[
\texttt{ship\_efield\_theory/src/propagation.py}
\]
函数名：
\[
\texttt{static\_dipole\_field(P\_a\_m, sigma\_s\_per\_m, R\_m, Ctheta)}
\]
函数实际做的计算就是：
\[
\texttt{return Ctheta * P\_a\_m / (4*pi*sigma\_s\_per\_m*R\_m**3)}
\]

主表生成工作流可以写成：
\[
\texttt{for R in [100,300,500,1000]:}
\]
\[
\texttt{\quad values=[]}
\]
\[
\texttt{\quad for P0 in P0\_list, sigma in sigma\_list, Ctheta in Ctheta\_list:}
\]
\[
\texttt{\quad\quad values.append(static\_dipole\_field(P0,sigma,R,Ctheta))}
\]
\[
\texttt{\quad min=min(values),\ reference=E(300,4,R,1),\ max=max(values)}
\]

\section{手算示例}
以 reference 场景、$R=100\ \mathrm{m}$ 为例：
\[
E=\frac{1\times300}{4\pi\times4\times100^3}
=5.97\times10^{-6}\ \mathrm{V/m}
\]
换成 $\mu$V/m：
\[
5.97\times10^{-6}\ \mathrm{V/m}\times10^6=5.97\ \mu\mathrm{V/m}
\]

\section{本源项输出如何进入主报告}
该源项进入表 5 的第一行。每个距离格写成：
\[
\mathrm{min/reference/max}
\]
其中 min 和 max 来自全部参数组合，reference 来自中等参考场景。

\end{document}
